\documentclass[10pt,a4paper]{article}
\usepackage{lmodern}
\usepackage[margin=19mm]{geometry}
\usepackage{booktabs}
\usepackage{graphicx}
\usepackage[table]{xcolor}
\usepackage[hidelinks]{hyperref}
\usepackage{parskip}
\usepackage{array}
\graphicspath{{figures/}}
\definecolor{warning}{HTML}{FFF3CD}
\definecolor{info}{HTML}{E0F2FE}
\hypersetup{pdftitle={Dynamic and static comparison of three UWB protocols},pdfauthor={UWB ESP-IDF technical report}}
\begin{document}
\begin{titlepage}
\centering
\vspace*{18mm}
{\Huge\bfseries Dynamic and static UWB comparison\par}
\vspace{4mm}
{\Large FlexTDOA, Native DS-TWR and Passive DS-TWR\par}
\vspace{14mm}
{\large Final field campaign of August 12, 2026\par}
\vspace{12mm}
\begin{tabular}{rl}
Platform: & 5 ESP32-S3 + DW3000 modules, Raspberry Pi collector\\
Topology: & M1 tag; M2, M3, M4 and M5 anchors\\
Radio: & UWB channel 9; configured 40 MHz SPI\\
Firmware branch: & \texttt{localization-raw-calibration}\\
UWB evidence: & six complete August 12 JSONL captures; raw and position-only EKF\\
RTK evidence: & GPS-only reference captures from August 6\\
\end{tabular}
\vfill
\colorbox{info}{\parbox{0.88\textwidth}{\vspace{2mm}
\textbf{Main result.} All three protocols produced complete raw dynamic and
static datasets. Independent dynamic rates are approximately 25.0 Hz Flex,
25.4 Hz Native and 22.6 Hz Passive. Raw static RMS is 2.3--2.4 cm. The
position-only adaptive EKF materially reduces static jitter and dynamic
sample-to-sample spikes, while raw UWB remains the permanent audit stream.\vspace{2mm}}}
\vfill
{\large Reproducible technical report\par}
\end{titlepage}
\pagenumbering{roman}
\tableofcontents
\clearpage
\pagenumbering{arabic}

\section{Executive summary}
\colorbox{warning}{\parbox{0.96\textwidth}{\textbf{Evidence boundary.}
Stream rates, uptime gaps, sample counts, fusion acceptance and reset counts are
direct August 12 diagnostics. GPS RTK quality and precision use only the August
6 reference captures. The two sessions are not time-paired, so this report does
not calculate or claim August 12 absolute UWB--RTK position error.}}

Passive DS overlapping raw windows remain visible in the telemetry rate, while
all comparative precision and trajectory metrics use independent frames.
Interleaved records marked \texttt{imu\_fused=true} are excluded from raw UWB
metrics. Outliers are retained.

The dynamic independent rates were \textbf{25.01 positions/s} for FlexTDOA, \textbf{25.35 positions/s} for Native DS-TWR, and \textbf{22.63 positions/s} for Passive DS-TWR. Passive DS emitted 70.11 dashboard events/s, but those include overlapping windows.

The static raw local RMS values were \textbf{2.41 cm} (FlexTDOA), \textbf{2.45 cm} (Native DS-TWR) and \textbf{2.26 cm} (Passive DS-TWR). These are precision around each capture median, not absolute accuracy.


The UWB and RTK evidence populations are intentionally separated: no August 6
UWB positions enter the protocol results, and no August 12 GPS samples enter
the RTK section. The new filter uses positions only: this campaign intentionally
does not claim acceleration-derived position improvement.

\section{Scope and questions}
The campaign asks how many independent positions each protocol delivers, where
stream gaps occur, how stable the unfiltered stationary position is, and whether
the position-only adaptive EKF reduces visible jitter without hiding raw UWB.

The RTK reference was acquired in a separate session. It establishes receiver
fix availability and short-term precision, but it cannot be paired sample by
sample with the August 12 UWB walks or used as an absolute-error ground truth.

\section{System and protocol background}
\subsection{FlexTDOA}
The tag is UWB receive-only. Anchor transmissions form directed time-difference
observations and the tag solves from a recent coherent observation set. Radio
airtime is largely independent of the number of passive tags.

\subsection{Native DS-TWR}
The active tag exchanges poll, response and final messages with every anchor,
obtains four absolute ranges and solves a coherent frame. Airtime grows with
active tag count, but frame semantics are direct.

\subsection{Passive DS-TWR}
The tag listens to anchor-to-anchor DS-TWR exchanges. Rolling windows may
produce several publications from shared measurements, so dashboard events and
independent frames are intentionally reported separately.

\section{Test configuration and captured evidence}
The three operator-marked dynamic walks lasted 85--114 s and the stationary
blocks approximately 30 s. The selector uses the final continuous-capture
segments and the final short static capture for each protocol.
\begin{center}\scriptsize\resizebox{\textwidth}{!}{%
\begin{tabular}{lrrrrrrrr}\toprule
\textbf{Protocol} & \textbf{Mode} & \textbf{Capture ID} & \textbf{duration s} & \textbf{events} & \textbf{independent} & \textbf{indep Hz} & \textbf{coverage pct} & \textbf{complete} \\
\midrule
FlexTDOA & dynamic & flextdoa\_final\_dynamic\_20260812 & 84.6 & 2117 & 2117 & 25.01 & 99.7 & yes \\
Native DS-TWR & dynamic & native\_ds\_final\_dynamic\_20260812 & 113.5 & 2876 & 2876 & 25.35 & 98.7 & yes \\
Passive DS-TWR & dynamic & passive\_ds\_final\_dynamic\_20260812 & 106.5 & 7463 & 2409 & 22.63 & 100.0 & yes \\
FlexTDOA & static & flextdoa\_static\_final\_01 & 30.3 & 798 & 798 & 26.30 & 98.8 & yes \\
Native DS-TWR & static & native\_ds\_static\_final\_02\_rtk\_fixed & 30.1 & 780 & 780 & 25.89 & 99.6 & yes \\
Passive DS-TWR & static & passive\_ds\_static\_final\_01 & 31.2 & 1906 & 607 & 19.45 & 96.2 & yes \\
\bottomrule\end{tabular}}\end{center}
Every selected block has a terminal \texttt{capture\_end}. All raw data are
archived losslessly and identified by SHA-256.

\section{Methodology}
\subsection{Population and continuity}
No error-based trimming or temporal smoothing is applied. Rates use full
capture wall time. Gaps are positive differences between unique ESP32 uptime
values; P50/P95/P99, maximum and counts above 100 ms are retained together.

\subsection{Static precision}
Radial displacement is computed around each block's component-wise position
median. CEP50, RMS, P95, P99, maximum, 2DRMS and first-to-last-decile drift use
all independent samples. They describe repeatability, not offset from truth.

\subsection{RTK reference}
Only quality-4 fixes from the August 6 reference captures enter the RTK section.
WGS84 is converted through ECEF to one local ENU frame for that session. No
cross-session timestamp association or UWB--RTK rigid fit is performed.

\section{Dynamic results}
\subsection{Raw trajectories}
The walks are separate physical trials, not synchronized repeats. The panels
therefore show continuity and gross geometry rather than point-by-point route
agreement.

\begin{figure}[htbp]
\centering
\includegraphics[width=\textwidth]{figures/06_dynamic_trajectories.pdf}
\caption{August 12 raw independent UWB trajectories in each local frame. Green markers are anchors; no August 6 UWB positions are used.}
\end{figure}
\subsection{Update rate and gaps}
\begin{center}\scriptsize\resizebox{\textwidth}{!}{%
\begin{tabular}{lrrrrrrrr}\toprule
\textbf{Protocol} & \textbf{duration s} & \textbf{events} & \textbf{independent} & \textbf{event Hz} & \textbf{indep Hz} & \textbf{gap P95 ms} & \textbf{gap max ms} & \textbf{step P95 m/s} \\
\midrule
FlexTDOA & 84.6 & 2117 & 2117 & 25.01 & 25.01 & 70 & 830 & 6.04 \\
Native DS-TWR & 113.5 & 2876 & 2876 & 25.35 & 25.35 & 50 & 130 & 7.15 \\
Passive DS-TWR & 106.5 & 7463 & 2409 & 70.11 & 22.63 & 30 & 250 & 4.66 \\
\bottomrule\end{tabular}}\end{center}
\begin{figure}[htbp]
\centering
\includegraphics[width=0.94\textwidth]{figures/01_position_rate.pdf}
\caption{Independent position rate in the selected dynamic and static blocks. Passive DS all-event publications are not substituted for independent frames.}
\end{figure}

\begin{figure}[htbp]
\centering
\includegraphics[width=0.94\textwidth]{figures/02_position_gap.pdf}
\caption{P95 positive gap between unique position uptimes. Maximum gaps and counts remain in the tables.}
\end{figure}

Native DS-TWR delivered the highest independent rate (25.35 Hz) and the best
worst-case continuity (130 ms). FlexTDOA was close in rate (25.01 Hz), but its
maximum gap was 830 ms. Passive DS delivered 22.63 independent Hz despite
overlapping dashboard publications, with a 250 ms maximum gap. Step speed is a
sample-to-sample trajectory diagnostic, not measured walking speed.

\subsection{Complete continuity audit}
\begin{center}\scriptsize\resizebox{\textwidth}{!}{%
\begin{tabular}{lrrrrrrrr}\toprule
\textbf{Protocol} & \textbf{Mode} & \textbf{event Hz} & \textbf{indep Hz} & \textbf{P50 ms} & \textbf{P95} & \textbf{P99} & \textbf{max} & \textbf{>100 ms} \\
\midrule
FlexTDOA & dynamic & 25.01 & 25.01 & 30 & 70 & 130 & 830 & 28 \\
Native DS-TWR & dynamic & 25.35 & 25.35 & 40 & 50 & 70 & 130 & 2 \\
Passive DS-TWR & dynamic & 70.11 & 22.63 & 10 & 30 & 80 & 250 & 28 \\
FlexTDOA & static & 26.30 & 26.30 & 30 & 70 & 100 & 120 & 7 \\
Native DS-TWR & static & 25.89 & 25.89 & 40 & 50 & 60 & 80 & 0 \\
Passive DS-TWR & static & 61.06 & 19.45 & 10 & 40 & 72 & 130 & 6 \\
\bottomrule\end{tabular}}\end{center}

\section{Position-only adaptive EKF}
The dashboard EKF uses only the raw UWB position stream. IMU acceleration is
disabled, raw measurements remain archived separately, and each EKF result in
this report comes from the same position event as its raw counterpart.
Sample-to-sample step-speed P95 is a jitter/continuity diagnostic; it is not
physical walking speed. Static RMS is repeatability about each stream's own
median. No cross-session RTK error is attached to the filter results.
\begin{center}\scriptsize\resizebox{\textwidth}{!}{%
\begin{tabular}{lrrrrrr}\toprule
\textbf{Protocol} & \textbf{Mode} & \textbf{coverage pct} & \textbf{raw step P95} & \textbf{EKF step P95} & \textbf{raw static RMS cm} & \textbf{EKF static RMS cm} \\
\midrule
FlexTDOA & dynamic & 100.0 & 6.04 & 4.31 & n/a & n/a \\
Native DS-TWR & dynamic & 100.0 & 7.15 & 3.44 & n/a & n/a \\
Passive DS-TWR & dynamic & 100.0 & 4.66 & 4.50 & n/a & n/a \\
FlexTDOA & static & 100.0 & 1.97 & 0.36 & 2.41 & 1.07 \\
Native DS-TWR & static & 100.0 & 1.61 & 0.23 & 2.45 & 0.89 \\
Passive DS-TWR & static & 100.0 & 1.62 & 0.00 & 2.26 & 0.00 \\
\bottomrule\end{tabular}}\end{center}
\begin{figure}[htbp]
\centering
\includegraphics[width=\textwidth]{figures/09_dynamic_raw_vs_ekf.pdf}
\caption{Raw independent dynamic trajectory (blue) and the position-only adaptive EKF (green). The EKF reduces high-frequency corners without replacing archived raw data.}
\end{figure}

\begin{figure}[htbp]
\centering
\includegraphics[width=\textwidth]{figures/10_static_raw_vs_ekf.pdf}
\caption{Static raw and EKF clouds, each centered on its own median and plotted at one common scale. Passive DS may collapse to a held point while stationary by design.}
\end{figure}


\section{Static results}
\subsection{Raw local precision}
\begin{center}\scriptsize\resizebox{\textwidth}{!}{%
\begin{tabular}{lrrrrrrrr}\toprule
\textbf{Protocol} & \textbf{N} & \textbf{CEP50 cm} & \textbf{RMS cm} & \textbf{P95 cm} & \textbf{P99 cm} & \textbf{max cm} & \textbf{2DRMS cm} & \textbf{drift cm} \\
\midrule
FlexTDOA & 798 & 1.91 & 2.41 & 4.30 & 5.73 & 7.38 & 4.81 & 0.82 \\
Native DS-TWR & 780 & 1.72 & 2.45 & 3.95 & 8.81 & 12.30 & 4.84 & 0.67 \\
Passive DS-TWR & 607 & 1.88 & 2.26 & 3.94 & 4.93 & 5.52 & 4.53 & 0.38 \\
\bottomrule\end{tabular}}\end{center}
\begin{figure}[htbp]
\centering
\includegraphics[width=0.94\textwidth]{figures/03_static_precision.pdf}
\caption{Static precision around each block's own median. All independent raw samples are included.}
\end{figure}

\begin{figure}[htbp]
\centering
\includegraphics[width=\textwidth]{figures/07_static_position_clouds.pdf}
\caption{Static raw position clouds centered on their own medians and drawn at one common scale. The cross is the median.}
\end{figure}

All raw static RMS values lie in a narrow 2.26--2.45 cm range. Passive DS is
best on RMS, while Native and Passive share the best P95 near 3.94 cm. All
outliers are retained. One stationary point cannot characterize GDOP or
spatial bias across the anchor polygon.

\section{GPS RTK reference from August 6}
This section contains GPS data only. It characterizes RTK availability and
short-term scatter for the unchanged receiver path; it does not reuse the old
UWB trajectories and it is not synchronized with the August 12 walks.
\begin{center}\scriptsize\resizebox{\textwidth}{!}{%
\begin{tabular}{lrrrrrrrrr}\toprule
\textbf{Protocol} & \textbf{Mode} & \textbf{tag total} & \textbf{tag fixed} & \textbf{tag fixed pct} & \textbf{all modules fixed pct} & \textbf{fix-age P95 ms} & \textbf{fixed anchors} & \textbf{tag cloud RMS cm} & \textbf{tag cloud P95 cm} \\
\midrule
FlexTDOA & dynamic & 24 & 24 & 100.0 & 93.0 & 96 & 4 & 284.65 & 482.34 \\
Native DS-TWR & dynamic & 30 & 28 & 93.3 & 55.9 & 96 & 3 & 334.51 & 487.56 \\
Passive DS-TWR & dynamic & 26 & 25 & 96.2 & 72.2 & 106 & 4 & 367.91 & 625.14 \\
FlexTDOA & static & 119 & 119 & 100.0 & 91.4 & 100 & 4 & 0.52 & 0.85 \\
Native DS-TWR & static & 119 & 119 & 100.0 & 86.3 & 100 & 4 & 0.74 & 1.15 \\
Passive DS-TWR & static & 119 & 102 & 85.7 & 93.0 & 130 & 4 & 1.40 & 2.02 \\
\bottomrule\end{tabular}}\end{center}
\begin{figure}[htbp]
\centering
\includegraphics[width=\textwidth]{figures/08_static_gps_rtk_clouds.pdf}
\caption{August 6 GPS RTK Fixed tag solutions only. Top: each reference block centered on its own median. Bottom: the same fixes in the August 6 common ENU frame.}
\end{figure}

The trajectory panels below are GPS-only and remain in the August 6 ENU frame.
They are shown as a qualitative RTK reference, not overlaid on the August 12
UWB paths.

\begin{figure}[htbp]
\centering
\includegraphics[width=\textwidth]{figures/11_dynamic_gps_rtk_reference.pdf}
\caption{August 6 GPS RTK Fixed dynamic reference trajectories. No UWB positions from that session are plotted.}
\end{figure}


\section{Accelerometer scope}
Acceleration-assisted positioning was disabled for this final campaign because
the preceding experiment did not improve raw UWB accuracy. The adaptive EKF
assessed above is position-only. Orientation/IMU telemetry is outside the
protocol ranking, and raw UWB remains the permanent audit stream.

\section{Protocol interpretation and selection}
\subsection{FlexTDOA}
Near-25 Hz independent dynamic rate with a receive-only scalable tag. Static
raw RMS was 2.41 cm. The 830 ms maximum gap is the principal caveat.

\subsection{Native DS-TWR}
Clearest coherent-frame semantics, highest measured independent rate and best
maximum gap. Static raw RMS was 2.45 cm. It remains the preferred deterministic
baseline for a small number of active tags.

\subsection{Passive DS-TWR}
Successfully follows motion without tag transmissions. Independent-frame
accounting yields 22.63 Hz despite overlapping publications; static raw RMS was
the best at 2.26 cm and the maximum dynamic gap was 250 ms.

\begin{center}\small
\begin{tabular}{p{0.20\textwidth}p{0.23\textwidth}p{0.23\textwidth}p{0.23\textwidth}}
\toprule
Criterion & FlexTDOA & Native DS-TWR & Passive DS-TWR\\
\midrule
Tag radio role & receive-only & active exchanges & receive-only\\
Dynamic independent rate & 25.01 Hz & best: 25.35 Hz & 22.63 Hz\\
Dynamic worst gap & 830 ms & best: 130 ms & 250 ms\\
Static raw precision & 2.41 cm RMS & 2.45 cm RMS & best: 2.26 cm RMS\\
Scaling with tag count & strongest & airtime grows & strong\\
Recommendation & primary scalable mode & coherent baseline & passive scalable alternative\\
\bottomrule
\end{tabular}
\end{center}

\section{Limitations}
\begin{itemize}
\item RTK is a separate August 6 reference session, not synchronized August 12 truth.
\item Cross-session UWB--RTK RMSE/P95 is intentionally not computed.
\item GNSS and UWB antenna phase centers have unknown lever arms.
\item Dynamic routes are not repeated surveyed paths; static/dynamic blocks are single trials.
\item One static point does not cover edge geometry, NLOS, orientation or spatial bias.
\item Collector rate is end-to-end telemetry, not direct DW3000 PHY throughput.
\end{itemize}

\section{Recommended next tests}
\begin{enumerate}
\item Repeat every protocol on one marked route with identical orientation and synchronized markers.
\item Interleave static blocks at the center, edges and outside the polygon.
\item Repeat the final UWB campaign with one simultaneous continuous RTK session.
\item Add hardware-correlated timestamps for latency claims.
\item Keep acceleration disabled until body-to-map yaw is independently validated.
\item Capture focused FlexTDOA diagnostics around gaps above 250 ms.
\end{enumerate}

\section{Preliminary findings}
\begin{itemize}
\item FlexTDOA remains a strong receive-only scalable mode, with one long-gap caveat.
\item Native DS-TWR is the best coherent-frame rate and continuity baseline.
\item Passive DS-TWR is a viable receive-only alternative when independent frames are counted correctly.
\item The cross-session RTK reference cannot support an August 12 absolute-accuracy ranking; local UWB precision is ranked.
\item The position-only EKF improves display stability; acceleration-assisted accuracy is not claimed.
\end{itemize}

\section{Reproduction and audit files}
The directory contains \texttt{analysis\_summary.json}, capture/replay/alignment
CSV files, SVG/PDF plots and this TeX source. The machine-readable summary
records whether optional lossless RAW archive generation was enabled; when it
is skipped, immutable source JSONL paths remain recorded for reproduction.

\begin{verbatim}
python tools/uwb_dynamic_static_report.py `
  --input-dir reports/uwb_final_report_input_20260812 `
  --output-dir reports/uwb_dynamic_static_comparison_20260812 `
  --rtk-reference-report-dir reports/uwb_dynamic_static_comparison_20260806
\end{verbatim}

No raw capture is modified during report generation.
\end{document}
